\acte{Les quinze oraisons de sainte Brigitte}
\usage{quot., pendant une année entière}
\begin{notice}
Selon la tradition, Notre-Seigneur enseigna ces quinze oraisons à sainte Brigitte, dans l’église Saint-Paul, à Rome, en mémoire des souffrances de sa Passion. Les « promesses » qui ont longtemps circulé avec elles ne sont pas reproduites ici : le Saint-Office a déclaré que leur origine surnaturelle n’est nullement établie et a demandé qu’on n’en permette plus la publication (monitum du 28 janvier 1954).
\end{notice}
\source{Texte : Les trois dévotions des prédestinés, 1888}
\rubrique{Avant chaque oraison, on dit un Pater et un Ave.}
\intertitre{Première oraison}
\ouverture{Ô}{}{Jésus-Christ !} douceur éternelle à ceux qui vous aiment, joie qui surpasse toute joie et tout désir, salut, ô cœur rempli d’amour pour les pécheurs, qui avez témoigné n’avoir point de plus grand contentement que d’être parmi les hommes, jusqu’à prendre la nature humaine dans la plénitude des temps pour leur amour. Souvenez-vous de la profonde tristesse que vous avez ressentie dès l’instant de votre sainte passion, ainsi qu’il avait été ordonné et décrété de toute éternité dans la pensée divine ; souvenez-vous de l’amertume qui a inondé votre âme, ainsi que vous le témoignâtes vous-même, lorsque faisant la Cène avec vos disciples, leur donnant votre précieux corps et sang, et leur lavant les pieds, leur prédites, sans d’autres consolations, vos prochaines souffrances. Souvenez-vous, mon Dieu, de toutes les craintes, angoisses et douleurs endurées dans votre humanité sainte avant votre mort, quand, après avoir prié par trois fois et après avoir sué sang et eau, vous fûtes trahi par Judas, votre disciple, pris par la nation que vous aviez choisie avec tant d’amour, accusé par de faux témoins, injustement condamné par Caïphe, Hérode et Pilate, en la fleur de votre jeunesse, et, dans le temps de Pâques, dépouillé de vos propres vêtements, et revêtu par dérision, de ceux d’autrui ; quand, les yeux et le visage couverts, vous reçûtes des soufflets ignominieux, vous fûtes couronné d’épines, attaché à une colonne ; quand on vous mit un roseau à la main, qu’on vous en frappa sur la tête, et qu’on vous accabla et qu’on vous déchira d’un nombre incalculable de calomnies et d’outrages. En mémoire de toutes ces peines et douleurs, endurées pour nous, dans votre passion, donnez-moi, avant la mort, une pure et entière confession, une vraie et sincère contrition, une digne satisfaction et la rémission de tous mes péchés. Ainsi soit-il.
\intertitre{Deuxième oraison}
\priere Ô Jésus ! vraie joie des Anges, Paradis de délices, ayez mémoire de l’horreur et de la tristesse que vous enduriez lorsque vos ennemis, tout ainsi que des lions furieux, vous entourèrent, et par injures, crachats, soufflets, égratignures, et mille autres mauvais traitements inouïs, vous tourmentèrent cruellement. Je vous prie, ô mon Seigneur, qu’en considération de tant de peines, paroles injurieuses et tourments, vous me délivriez de tous mes ennemis, tant visibles qu’invisibles, et me fassiez la grâce qu’étant sous votre protection, je puisse parvenir à la possession du salut éternel. Ainsi soit-il.
\intertitre{Troisième oraison}
\priere Ô Jésus ! créateur du ciel et de la terre, que nulle chose ne peut borner ni limiter ; vous qui renfermez tout dans votre puissance, souvenez-vous de la douleur très vive que vous ressentîtes lorsque les Juifs attachant vos mains sacrées et vos pieds très délicats à la Croix, les percèrent de part en part, avec de gros clous émoussés, parce que vous condamniez leurs volontés charnelles ; et qu’ils ajoutèrent à vos plaies douleurs sur douleurs, et en cette façon vous ont si cruellement étendu sur la Croix, que toutes les jointures de vos membres en furent séparées. Je vous conjure, par la mémoire de cette très amère douleur de la Croix, de me donner votre crainte et votre amour. Ainsi soit-il.
\intertitre{Quatrième oraison}
\priere Ô Jésus ! céleste médecin, souvenez-vous des langueurs et meurtrissures souffertes en tous vos membres, dont pas un ne demeura en sa place quand vous fûtes élevé en Croix, de façon qu’il n’y eut douleur semblable à la vôtre, d’autant qu’il ne s’est rencontré en vous, depuis la plante des pieds jusqu’au sommet de la tête, aucune partie de votre corps qui n’ait enduré ; et cependant oubliant toutes vos douleurs, vous n’avez pas laissé de prier votre Père pour vos ennemis, lui disant : Mon Père, pardonnez-leur, car ils ne savent ce qu’ils font. Par cette grande miséricorde, et en mémoire de cette douleur, faites, je vous en supplie, que le souvenir de votre passion très douloureuse me procure une pleine rémission de tous mes péchés. Ainsi soit-il.
\intertitre{Cinquième oraison}
\priere Ô Jésus ! miroir de la lumière éternelle, souvenez-vous de la tristesse que vous avez eue, lorsque vous regardâtes et contemplâtes dans le miroir de votre divinité, la prédestination de ceux qui doivent être sauvés par les mérites de votre sainte Passion, et la grande multitude des réprouvés qui doivent être damnés pour leurs péchés : par cet abîme de compassion et de pitié, dont vous avez regretté et plaint ces malheureux pécheurs perdus et désespérés, et principalement par la bonté que vous fîtes paraître en Croix au bon larron, en lui disant : Tu seras aujourd’hui avec moi en Paradis, je vous prie, ô doux Jésus, d’user, à l’heure de ma mort, de la même miséricorde envers moi. Ainsi soit-il.
\intertitre{Sixième oraison}
\priere Ô Jésus ! roi aimable, souvenez-vous de la douleur que vous avez eue quand dépouillé de vos vêtements et comme un malfaiteur, fûtes attaché à la Croix, où vos proches et amis vous abandonnèrent et demeurâtes sans consolation, sinon de quelques dévouées et de votre Mère bien-aimée que vous recommandâtes à votre disciple saint Jean, en lui disant : Femme, voilà votre Fils. Je vous supplie, ô mon Sauveur, par le glaive de douleur qui pour lors transperça l’âme de cette douce Mère, que vous ayez compassion de moi, et que vous me soulagiez en toutes mes tribulations et afflictions, tant corporelles que spirituelles, et me donniez consolation dans le temps de mes tourments et à l’heure de ma mort. Ainsi soit-il.
\intertitre{Septième oraison}
\priere Ô Jésus ! fontaine de pitié inépuisable, qui par une profonde affection d’amour avez dit en Croix : j’ai soif, mais de la soif du salut du genre humain, je vous prie, ô mon Sauveur, d’échauffer le désir de nos cœurs, de telle sorte que toutes nos œuvres et actions soient parfaites, et d’éteindre entièrement en nous la concupiscence charnelle, et de refroidir l’ardeur des appétits mondains. Ainsi soit-il.
\intertitre{Huitième oraison}
\priere Ô Jésus ! douceur des cœurs, suavité des esprits, par l’amertume du fiel et du vinaigre que vous avez goûtés sur la Croix pour l’amour de nous, permettez, je vous en supplie, que nous puissions recevoir dignement votre corps et votre sang précieux à l’heure de notre mort, pour qu’ils servent de remède et de consolation à nos âmes. Ainsi soit-il.
\intertitre{Neuvième oraison}
\priere Ô Jésus ! vertu royale, joie de l’esprit, souvenez-vous de l’angoisse et de la douleur que vous avez endurées, lorsque, par l’ignominie de votre mort, par les moqueries et les outrages des Juifs, vous criâtes à haute voix que vous aviez été abandonné de votre Père, lui disant : Mon Dieu, mon Dieu, pourquoi m’avez-vous délaissé ? Par cette tristesse et cette amertume, je vous conjure, ô mon Sauveur, de ne me point délaisser dans toutes les angoisses et douleurs de ma mort. Ainsi soit-il.
\intertitre{Dixième oraison}
\priere Ô Jésus ! qui êtes en toutes choses commencement et fin, vie et vertu, souvenez-vous que vous avez été abîmé et submergé dans l’eau de votre Passion, depuis le sommet de la tête jusqu’à la plante des pieds. En considération de la grandeur de vos plaies, enseignez-moi à garder vos commandements dans une vraie charité, afin que je ne sois point submergé et abîmé par mes péchés. Ainsi soit-il.
\intertitre{Onzième oraison}
\priere Ô Jésus ! abîme très profond de miséricorde, je vous prie qu’en mémoire de vos plaies, de ces plaies dont les douleurs ont passé jusqu’à la moelle de vos os et jusqu’au fond de vos entrailles, de ne point permettre que je sois submergé, misérable que je suis, dans le gouffre de mes offenses ; mais retirez-moi, au contraire, hors du péché, et cachez-moi loin de votre face irritée dans vos saintes plaies, jusqu’à ce que votre colère et votre indignation soient passées. Ainsi soit-il.
\intertitre{Douzième oraison}
\priere Ô Jésus ! miroir de vérité, marque d’unité, lien de charité, souvenez-vous de la multitude des plaies dont vous avez été blessé depuis les pieds jusqu’à la tête, quand vous avez été déchiré et baigné dans les flots de votre sang adorable et que vous avez souffert pour l’amour de nous en votre chair virginale. Ô doux Jésus, que pouviez-vous faire pour nous, que vous n’ayez fait ? Je vous conjure, ô mon Sauveur, de marquer avec votre précieux sang toutes vos plaies dans mon cœur, afin que par là je connaisse clairement vos douleurs et votre amour, et qu’en action de grâces je persévère courageusement jusqu’à la fin. Ainsi soit-il.
\intertitre{Treizième oraison}
\priere Ô Jésus ! Dieu très fort, roi immortel et invincible, souvenez-vous de la douleur que vous avez endurée, lorsque toutes vos forces entièrement affaiblies, vous dites en inclinant la tête : Tout est consommé. Par cette angoisse et cette douleur, je vous prie, ô mon bon Jésus, d’avoir pitié de moi en la dernière heure de ma vie, lorsque mon âme sera dans l’angoisse et que mon esprit sera troublé. Ainsi soit-il.
\intertitre{Quatorzième oraison}
\priere Ô Jésus ! fils unique du Père, la splendeur et la figure de sa substance, souvenez-vous de l’humble recommandation que vous fîtes à votre Père, quand vous dîtes : Je remets mon esprit entre vos mains ; et puis, le corps déchiré, le cœur brisé et les entrailles de votre miséricorde ouvertes pour nous racheter, vous avez expiré. Par cette précieuse mort, je vous prie, ô Roi des saints, fortifiez-moi et donnez-moi secours pour résister au démon, au monde, à la chair et sang, afin qu’étant mort au monde, je vive à vous seul. Recevez, je vous prie, à l’heure de ma mort, mon âme pèlerine et exilée, qui veut retourner à vous. Ainsi soit-il.
\intertitre{Quinzième oraison}
\priere Ô Jésus ! vraie et féconde vigne, souvenez-vous de la grande effusion de sang que vous avez répandu de votre corps adorable, tout ainsi que le raisin mis sous le pressoir ; et lorsque votre côté fut percé d’un coup de lance par un soldat, vous avez donné sang et eau, en telle sorte qu’il n’en est pas demeuré une seule goutte dans votre cœur ; ainsi, comme un faisceau de myrrhe, vous avez été élevé au haut de la Croix, votre chair délicate est devenue à néant, l’humeur de vos entrailles s’est toute tarie, et la moelle de vos os s’est séchée. Par ces amers traitements et par l’effusion de votre précieux sang, ô mon bon Jésus ! je vous prie de recevoir mon âme lorsque je serai à l’agonie de la mort. Ainsi soit-il.
\intertitre{Conclusion}
\priere Ô doux Jésus ! blessez mon cœur, afin que mes larmes de pénitence, de douleur et d’amour, nuit et jour me servent de pain ; convertissez-moi entièrement à vous ; que mon cœur vous soit une perpétuelle habitation ; que ma conversation vous soit agréable, et que la fin de ma vie vous soit tellement louable, qu’après ma mort je puisse mériter votre Paradis, et vous louer à jamais dans l’éternité. Ainsi soit-il.
\culdelampe